
\documentclass[a4paper,12pt]{article}
\usepackage[czech,slovak]{babel}
\usepackage[a4paper, total={170mm, 240mm}, top=2cm, left=3cm, right=3cm]{geometry}
\usepackage[utf8]{inputenc}
\usepackage[IL2]{fontenc}
\usepackage{times}
\usepackage[hidelinks]{hyperref}
\usepackage{graphicx}
\usepackage{booktabs}
\usepackage{epstopdf}
\usepackage[htt]{hyphenat}
\usepackage{amsmath}
\usepackage{float}
\usepackage{indentfirst}
\usepackage{array}
\usepackage{enumitem}
\usepackage{bm}

\begin{document}
\begin{center}
\includegraphics[width=1.0\textwidth]{fit.pdf} % змініть розмір при необхідності

    \vspace{0.5cm} 
\Huge
\textsc{Vysoké učení technické v~Brně\\
}Fakulta informačních technologií\\
\vspace{\stretch{0.382}}
\Huge Technická správa \\
\LARGE T5: Zábavní průmysl \\
\LARGE  Model parku Walt Disney Studio v Disneylandu v Paříži\\
\vspace{\stretch{0.309}}

\Large 

\vspace{\stretch{0.309}}

\end{center}
{\Large \today \hfill
\begin{tabular}{l}
    Maryna Kucher, xkuche01 \\
    Vladyslava Bilyk, xbilyk03
\end{tabular}}
\thispagestyle{empty}

\newpage
\tableofcontents
\newpage

\section{Úvod}
V rámci tohoto projektu byla vytvořena modelová simulace jednoho z parků v Disneylandu v Paříži. Park nese název Walt Disney Studio a obsahuje 10 atrakcí a N tematických sekcí. Náš model ukazuje vstup do parku, volbu atrakce, čekání ve frontě na atrakci a samotnou jízdu na ní.

Cílem naší simulace bylo porovnat zisk, který park generuje od návštěvníků za současného nastavení (tj. cena za vstupenku na celý park), s variantou, kdy návštěvníci platí za vstup na každou atrakci zvlášť. Smyslem experimentů je demonstrovat, že pokud by Disneyland vybíral poplatek za každou atrakci zvlášť, místo jednorázového vstupného do samotného parku, dosáhl by vyššího zisku.

\subsection{Autoři a zdroje}
Tuto práci zpracovali Vladyslava Bilyk a Maryna Kucher. Informace byly získány z osobní návštěvy parku a byly podpořeny statistickými údaji z článků a internetových stránek.

\subsection{Validace}
Validace našeho programu byla provedena následujícím způsobem: experimentálně, po několika spuštěních modelu, byly statistické výstupy dat porovnány s očekávanými údaji, které odpovídají realitě. Každá fáze modelování byla analyzována z hlediska souladu s reálnými procesy. Vstupní parametry pro simulaci, jako je cena za vstupenku, počet návštěvníků za den, počet atrakcí, jejich kapacita a některé další údaje, byly čerpány z oficiálních stránek Disneylandu.

\section{Rozbor tématu a použitých metod/technologií}

\section{Koncepce}
\subsection{Koncepce - modelářská témata}
\subsection{Koncepce - implementační témata}

\section{Architektura simulačního modelu/simulátoru}

\section{Podstata simulačních experimentů a jejich průběh}

\section{Shrnutí simulačních experimentů a závěr}
\bibliographystyle{czechiso}
\bibliography{citations.bib} 

\end{document}
